\documentclass[12pt,a4paper]{article}
\usepackage[utf8]{inputenc}
\usepackage[T1]{fontenc}
\usepackage[hungarian]{babel}
\usepackage{geometry}
\geometry{margin=2.5cm}
\usepackage{graphicx}
\usepackage{booktabs}
\usepackage{hyperref}
\usepackage{bookmark}
\usepackage{xcolor}
\usepackage{listings}
\usepackage{amsmath}
\usepackage{float}
\usepackage{enumitem}

\hypersetup{
    colorlinks=true,
    linkcolor=blue!60!black,
    urlcolor=blue!60!black,
    citecolor=blue!60!black,
}

\lstset{
    basicstyle=\ttfamily\small,
    keywordstyle=\color{blue!70!black}\bfseries,
    commentstyle=\color{green!50!black}\itshape,
    stringstyle=\color{red!60!black},
    backgroundcolor=\color{gray!8},
    frame=single,
    framerule=0.4pt,
    rulecolor=\color{gray!50},
    breaklines=true,
    showstringspaces=false,
    numbers=left,
    numberstyle=\tiny\color{gray},
    language=Python,
    tabsize=4,
}

\title{
    \textbf{Ipari Képfeldolgozás Laboratóriumi Gyakorlat} \\[0.3cm]
    \Large 2.~Mérföldkő — Kezdeti Implementáció \& Első Eredmények \\[0.3cm]
    \large Projekt: Objektumok regisztrációja 3D pontfelhőn
}
\author{}
\date{}

\begin{document}
\maketitle
\thispagestyle{empty}

\vspace{-1.5cm}

\begin{table}[H]
\centering
\begin{tabular}{ll}
\toprule
\textbf{Beadási határidő} & 2026.~április 15.~23:59 \\
\textbf{Csapat tagjai}     & Ilonczai András \\
\textbf{Dátum}             & 2026.~április \\
\bottomrule
\end{tabular}
\end{table}

% ============================================================
\section{Projekt összefoglalás és státusz}
% ============================================================

Az 1.~mérföldkőben bemutatott algoritmus terv és költségbecslés alapján elkészült a teljes 3D pontfelhő regisztrációs pipeline kezdeti implementációja. A rendszer Python nyelven, az Open3D könyvtárra építve valósítja meg a Stanford Bunny adathalmaz 10 különböző szögből felvett pontfelhőjének automatikus illesztését és egyetlen koherens 3D modellé történő összefűzését.

A jelenlegi állapotban a pipeline az alábbi funkcionalitásokat tartalmazza:
\begin{itemize}[nosep]
    \item Teljes feldolgozási lánc (betöltéstől az exportálásig)
    \item RANSAC durva illesztés + ICP finom illesztés (Point-to-Point és Point-to-Plane)
    \item Automatikus minőségértékelés (fitness, RMSE)
    \item Vizualizáció (Matplotlib 2D + PyVista 3D off-screen renderelés)
    \item Interaktív 3D megjelenítő (Open3D visualizer)
\end{itemize}

% ============================================================
\section{Implementált algoritmus}
% ============================================================

\subsection{Pipeline architektúra}

A megvalósított pipeline az 1.~mérföldkő tervét követi, hat fő lépésben:

\begin{table}[H]
\centering
\begin{tabular}{clp{8cm}}
\toprule
\textbf{Lépés} & \textbf{Megnevezés} & \textbf{Leírás} \\
\midrule
1 & Adatbevitel \& előfeldolgozás & PLY fájlok beolvasása, voxel grid ritkítás ($v = 0{,}005$~m), Statistical Outlier Removal \\
2 & Jellemzők kinyerése & PCA-alapú normális becslés, FPFH deszkriptorok (33D) \\
3 & Durva illesztés & RANSAC globális regisztráció FPFH jellemzőkkel \\
4 & Finom illesztés & ICP Point-to-Point és Point-to-Plane variánsok \\
5 & Minőségértékelés & Fitness (overlap arány), inlier RMSE, korespondenciák száma \\
6 & Eredmény exportálás & Merged PLY fájl, transzformációs mátrixok (JSON), log \\
\bottomrule
\end{tabular}
\caption{A pipeline hat fő feldolgozási lépése.}
\end{table}

\subsection{Adatbevitel és előfeldolgozás (1.~lépés)}

A Stanford Bunny adathalmaz egy \texttt{bun.conf} konfigurációs fájlt tartalmaz, amely 10 szkennt definiál ground-truth transzformációkkal (pozíció + quaternion orientáció).

\paragraph{Konfigurációs fájl formátuma:}
\begin{lstlisting}[language={}]
bmesh <fájlnév> tx ty tz  qx qy qz qw
\end{lstlisting}

Az egyes szkenneket az alábbi előfeldolgozási lépéseken vezetjük át:

\begin{enumerate}[nosep]
    \item \textbf{Beolvasás:} Open3D \texttt{read\_point\_cloud()} PLY fájlok betöltéséhez.
    \item \textbf{Voxel grid ritkítás:} $v = 0{,}005$~m (5~mm) méretű voxelekkel, amely jelentősen csökkenti a pontfelhő méretét a számítási hatékonyság érdekében.
    \item \textbf{Statistical Outlier Removal (SOR):} 20 szomszéd, $2\sigma$ küszöbérték. Eltávolítja a zajos, izolált pontokat.
\end{enumerate}

\subsection{Jellemzők kinyerése (2.~lépés)}

\paragraph{Normális becslés:}
PCA-alapú normális becslés hibrid keresési paraméterekkel:
\begin{itemize}[nosep]
    \item Keresési sugár: $r = 2v = 0{,}01$~m
    \item Maximális szomszédok: 30
    \item Konzisztens orientáció: \texttt{orient\_normals\_consistent\_tangent\_plane(k=30)}
\end{itemize}

\paragraph{FPFH deszkriptorok:}
A Fast Point Feature Histograms (FPFH) \cite{rusu2009} 33 dimenziós hisztogramban kódolja a lokális geometriai jellemzőket:
\begin{itemize}[nosep]
    \item Keresési sugár: $r = 5v = 0{,}025$~m
    \item Maximális szomszédok: 100
\end{itemize}

Az FPFH robusztus jellemző-egyeztetést tesz lehetővé különböző nézőpontok között.

\subsection{Durva illesztés — RANSAC (3.~lépés)}

A RANSAC (Random Sample Consensus) \cite{zhou2016} algoritmus az FPFH jellemzők közötti megfeleltetésekből becsüli a kezdeti merev transzformációt:

\begin{table}[H]
\centering
\begin{tabular}{ll}
\toprule
\textbf{Paraméter} & \textbf{Érték} \\
\midrule
Távolság küszöb & $1{,}5v = 0{,}0075$~m \\
Maximális iteráció & 100\,000 \\
Konfidencia & 0,999 \\
Minta pontok (ransac\_n) & 3 \\
\bottomrule
\end{tabular}
\caption{RANSAC paraméterek.}
\end{table}

Az ellenőrző feltételek (\textit{checkers}):
\begin{itemize}[nosep]
    \item \textbf{Edge length:} A három mintapont közötti élek arányainak konzisztenciája ($\geq 0{,}9$).
    \item \textbf{Distance:} A transzformált pontpárok közötti távolság ellenőrzése.
\end{itemize}

\subsection{Finom illesztés — ICP (4.~lépés)}

Az Iterative Closest Point (ICP) \cite{besl1992} algoritmus a RANSAC eredményéből kiindulva minimalizálja a forrás- és célpontfelhő közötti távolságot.

Két variánst implementáltunk:

\paragraph{Point-to-Point ICP:}
\begin{equation}
    E(\mathbf{R}, \mathbf{t}) = \sum_{i=1}^{N} \| \mathbf{R} \mathbf{p}_i + \mathbf{t} - \mathbf{q}_i \|^2
\end{equation}

\paragraph{Point-to-Plane ICP:}
\begin{equation}
    E(\mathbf{R}, \mathbf{t}) = \sum_{i=1}^{N} \left[ (\mathbf{R} \mathbf{p}_i + \mathbf{t} - \mathbf{q}_i) \cdot \mathbf{n}_i \right]^2
\end{equation}

ahol $\mathbf{p}_i$ a forrás pontok, $\mathbf{q}_i$ a legközelebbi célpontok, és $\mathbf{n}_i$ a célpont normálisai \cite{chen1992}.

\begin{table}[H]
\centering
\begin{tabular}{ll}
\toprule
\textbf{Paraméter} & \textbf{Érték} \\
\midrule
Távolság küszöb & $0{,}4v = 0{,}002$~m \\
Maximális iteráció & 200 \\
Relatív fitness konvergencia & $10^{-6}$ \\
Relatív RMSE konvergencia & $10^{-6}$ \\
\bottomrule
\end{tabular}
\caption{ICP paraméterek.}
\end{table}

\subsection{Multi-scan rekonstrukció}

A teljes rekonstrukció inkrementálisan történik:
\begin{enumerate}[nosep]
    \item A \texttt{bun000.ply} a referencia (target) scan.
    \item Minden további scan a ground-truth transzformáció alkalmazásával kap kezdő pozíciót.
    \item ICP Point-to-Plane finomítja az illesztést a már meglévő fúzionált felhőhöz.
    \item A legjobb transzformáció kiválasztása (GT vs.\ ICP) a fitness alapján.
    \item A transzformált scan hozzáfűzése a merged felhőhöz, voxel ritkítással.
\end{enumerate}

\subsection{Minőségértékelés (5.~lépés)}

Három metrikát számolunk minden illesztéshez:

\begin{itemize}[nosep]
    \item \textbf{Fitness (overlap arány):} Az inlier megfeleltetések aránya a forrás pontszámához képest. Ideális: $\geq 0{,}8$.
    \item \textbf{Inlier RMSE:} Az inlier pontpárok közötti négyzetes középhiba (méterben). Ideális: $< 0{,}002$~m.
    \item \textbf{Korespondenciák száma:} A távolság-küszöbön belüli pontpárok száma.
\end{itemize}

% ============================================================
\section{Implementáció részletei}
% ============================================================

\subsection{Projekt struktúra}

\begin{lstlisting}[language={}]
2. merfoldko/
  registration_pipeline.py   # Fo pipeline (~630 sor)
  visualize.py               # Interaktiv 3D megjelenito
  setup_env.bat              # Kornyezet telepitese (Windows)
  requirements.txt           # Python fuggosegek
  commands.txt               # Hasznos parancsok
  input/
    bunny/
      data/
        bun.conf             # 10 scan konfiguracioja
        bun000.ply ... ear_back.ply  # PLY pontfelhok
  output/                    # Generalt eredmenyek (gitignore-ban)
\end{lstlisting}

\subsection{Felhasznált könyvtárak}

\begin{table}[H]
\centering
\begin{tabular}{lll}
\toprule
\textbf{Könyvtár} & \textbf{Verzió} & \textbf{Felhasználás} \\
\midrule
Open3D & $\geq$0.18.0 & Pontfelhő I/O, ICP, RANSAC, vizualizáció \\
NumPy & $\geq$1.24.0 & Numerikus számítások, mátrixműveletek \\
SciPy & $\geq$1.11.0 & Quaternion konverzió (\texttt{Rotation}) \\
Matplotlib & $\geq$3.7.0 & 2D ábrák, minőségdiagramok \\
PyVista & $\geq$0.43.0 & Off-screen 3D renderelés \\
\bottomrule
\end{tabular}
\caption{Python függőségek.}
\end{table}

\subsection{Főbb implementációs döntések}

\begin{enumerate}
    \item \textbf{Voxel méret ($v = 5$~mm):} A Stanford Bunny ~0,2~m magas. Az 5~mm-es voxel jó egyensúlyt ad a részletesség és a számítási sebesség között.
    \item \textbf{Ground-truth inicializáció:} A multi-scan fúzió a \texttt{bun.conf} fájl GT transzformációit használja kezdő becslésként, amit ICP finomít. Ez robusztusabb, mint a puszta RANSAC, mivel a Bunny szimmetriája félrevezetheti a globális regisztrációt.
    \item \textbf{Inkrementális fúzió:} Minden scant az eddigi merged felhőhöz illesztünk, nem pedig páronként. Ez csökkenti a kumulatív hibát.
    \item \textbf{Automatikus minőségválasztás:} Ha az ICP nem javít a GT-n (fitness alapján), a rendszer a GT transzformációt tartja meg.
\end{enumerate}

% ============================================================
\section{Bemeneti adathalmaz}
% ============================================================

A Stanford Bunny adathalmaz \cite{open3d2018} 10 különböző nézőpontból tartalmaz pontfelhőket:

\begin{table}[H]
\centering
\begin{tabular}{ll}
\toprule
\textbf{Fájl} & \textbf{Leírás} \\
\midrule
\texttt{bun000.ply} & Referencia nézet (0°) \\
\texttt{bun045.ply} & 45° elforgatás \\
\texttt{bun090.ply} & 90° elforgatás \\
\texttt{bun180.ply} & 180° elforgatás \\
\texttt{bun270.ply} & 270° elforgatás \\
\texttt{bun315.ply} & 315° elforgatás \\
\texttt{top2.ply}   & Felülnézet 2 \\
\texttt{top3.ply}   & Felülnézet 3 \\
\texttt{chin.ply}   & Áll alulnézete \\
\texttt{ear\_back.ply} & Fül hátsó oldala \\
\bottomrule
\end{tabular}
\caption{A Stanford Bunny adathalmaz 10 scan fájlja.}
\end{table}

% ============================================================
\section{Első eredmények}
% ============================================================

\subsection{Pipeline futtatás}

A pipeline futtatása:
\begin{lstlisting}[language=bash]
cd "2. merfoldko"
.\setup_env.bat
.venv\Scripts\python.exe registration_pipeline.py
\end{lstlisting}

\subsection{Generált eredmények}

A pipeline az \texttt{output/} könyvtárba az alábbi fájlokat generálta:

\begin{itemize}[nosep]
    \item \texttt{bunny\_full\_reconstruction.ply} — a 10 scan összefűzött pontfelhője (2\,585 pont)
    \item \texttt{transforms.json} — végső transzformációs mátrixok (fájlnév → 4$\times$4 mátrix)
    \item \texttt{pipeline\_log.txt} — részletes futási napló (időbélyeges)
\end{itemize}

\subsection{Előfeldolgozási statisztikák}

Az egyes szkennelések pontszámának alakulása az előfeldolgozás során:

\begin{table}[H]
\centering
\begin{tabular}{lccc}
\toprule
\textbf{Scan} & \textbf{Eredeti} & \textbf{Voxel után} & \textbf{SOR után} \\
\midrule
bun000.ply     & 40\,256 & 1\,408 & 1\,348 \\
bun045.ply     & 40\,097 & 1\,335 & 1\,269 \\
bun090.ply     & 30\,379 & 1\,251 & 1\,217 \\
bun180.ply     & 40\,251 & 1\,365 & 1\,342 \\
bun270.ply     & 31\,701 & 1\,243 & 1\,194 \\
top2.ply       & 38\,298 & 1\,397 & 1\,337 \\
top3.ply       & 36\,023 & 1\,275 & 1\,257 \\
bun315.ply     & 35\,336 & 1\,387 & 1\,348 \\
chin.ply       & 37\,738 & 1\,425 & 1\,370 \\
ear\_back.ply  & 32\,193 & 1\,173 & 1\,168 \\
\midrule
\textbf{Összesen} & \textbf{362\,272} & \textbf{13\,259} & \textbf{12\,850} \\
\bottomrule
\end{tabular}
\caption{Pontszám változása az előfeldolgozás során. A voxel ritkítás ($v = 5$~mm) és SOR szűrés a pontszámot átlagosan 96,5\%-kal csökkentette.}
\end{table}

\subsection{Regisztrációs minőség}

Az alábbi táblázat a pipeline futtatásának tényleges eredményeit tartalmazza. Minden scant a már meglévő fúzionált felhőhöz illesztettünk; a rendszer automatikusan a jobb fitness-értékű módszert választotta (ICP P2Plane refinement vs.\ puszta GT).

\begin{table}[H]
\centering
\begin{tabular}{lcccc}
\toprule
\textbf{Scan} & \textbf{Fitness} & \textbf{Inlier RMSE (m)} & \textbf{\#Corr} & \textbf{Módszer} \\
\midrule
bun045.ply     & 0,9125 & 0,001851 & 1\,158 & ICP P2Plane \\
bun090.ply     & 0,6122 & 0,002033 &   745 & ICP P2Plane \\
bun180.ply     & 0,4858 & 0,002442 &   652 & ICP P2Plane \\
bun270.ply     & 0,9497 & 0,002320 & 1\,134 & GT \\
bun315.ply     & 0,9918 & 0,002062 & 1\,337 & ICP P2Plane \\
top2.ply       & 0,9783 & 0,002327 & 1\,308 & ICP P2Plane \\
top3.ply       & 1,0000 & 0,002227 & 1\,257 & ICP P2Plane \\
chin.ply       & 0,7686 & 0,002224 & 1\,053 & ICP P2Plane \\
ear\_back.ply  & 1,0000 & 0,002185 & 1\,168 & ICP P2Plane \\
\midrule
\textbf{Átlag} & \textbf{0,8554} & \textbf{0,002186} & & \\
\bottomrule
\end{tabular}
\caption{Multi-scan regisztrációs eredmények (pipeline futtatás: 2026-04-14).}
\end{table}

\paragraph{Eredmények elemzése:}
\begin{itemize}[nosep]
    \item A legjobb illesztést a \texttt{top3.ply} és \texttt{ear\_back.ply} érte el (fitness = 1,0000), ami teljes átfedést jelez a meglévő modellel.
    \item A \texttt{bun315.ply} és \texttt{bun270.ply} szintén kiváló eredményt mutat (fitness $>$ 0,94).
    \item A \texttt{bun180.ply} érte el a legalacsonyabb fitness-értéket (0,4858), ami a 180°-os nézet korlátozott átfedésére utal az addigi modellel.
    \item Az ICP P2Plane refinement 8 esetből 8-ban javított vagy azonos eredményt adott a GT-hez képest; egyedül a \texttt{bun270.ply} esetében maradt a GT a jobb választás.
    \item Az átlagos inlier RMSE 0,002186~m ($\approx$2,2~mm), ami az 5~mm-es voxel mérethez képest jó pontosságot jelent.
    \item A végső rekonstrukció \textbf{2\,585 pontot} tartalmaz (voxel-ritkítás után).
\end{itemize}

\subsection{Vizualizáció}

Az interaktív megjelenítő futtatása:
\begin{lstlisting}[language=bash]
.venv\Scripts\python.exe visualize.py
\end{lstlisting}

A megjelenítő lehetőségei:
\begin{itemize}[nosep]
    \item Egyedi input szkennelések megtekintése (magasság szerinti színezéssel)
    \item A teljes rekonstrukció megjelenítése Poisson felületrekonstrukcióval
    \item Interaktív forgatás, zoom, pont méretezés
\end{itemize}

% ============================================================
\section{Értékelés és következő lépések}
% ============================================================

\subsection{Jelenlegi eredmények értékelése}

\begin{itemize}
    \item A pipeline sikeresen beolvassa mind a 10 Stanford Bunny scant.
    \item Az előfeldolgozás (voxel + SOR) hatékonyan csökkenti a pontszámot és a zajt.
    \item A RANSAC + ICP kombináció megbízhatóan illeszti a szomszédos nézeteket.
    \item A ground-truth inicializáció biztosítja a robusztus konvergenciát, különösen a top és ear\_back nézeteknél, amelyeknél a geometriai átfedés korlátozott.
\end{itemize}

\subsection{Ismert korlátok}

\begin{itemize}
    \item A szimmetrikus részek (pl.\ a nyúl két oldala) félrevezethetik a puszta RANSAC-ot GT inicializáció nélkül.
    \item Az inkrementális fúzió kumulatív hibát halmozhat fel sok scan esetén (drift).
    \item A voxel-ritkítás finom geometriai részleteket (fül, orr) elmoshat.
\end{itemize}

\subsection{3.~mérföldkőre tervezett fejlesztések}

\begin{enumerate}[nosep]
    \item Globális optimalizáció (pose graph) a drift csökkentéséhez.
    \item Colored ICP kipróbálása (ha textúra-információ elérhető).
    \item Több adathalmaz tesztelése (nem csak Stanford Bunny).
    \item Paraméter-érzékenységi elemzés (voxel méret, ICP küszöb hatása).
    \item A végleges eredmények számszerű összefoglalása és összehasonlítása.
\end{enumerate}

% ============================================================
\section{Projekt weboldal}
% ============================================================

A projekt haladása egy dedikált weboldalon követhető, mérföldkőnként elkülönítve:

\begin{itemize}[nosep]
    \item \textbf{1.~mérföldkő:} Algoritmus terv és költségbecslés
    \item \textbf{2.~mérföldkő:} Kezdeti implementáció és első eredmények (jelen dokumentum)
    \item \textbf{3.~mérföldkő:} Végleges eredmények és értékelés
\end{itemize}

\textbf{Weboldal URL:} \textit{(kitöltendő)}

% ============================================================
\begin{thebibliography}{9}
\bibitem{besl1992}
Besl, P.~J. \& McKay, N.~D. (1992).
A method for registration of 3-D shapes.
\textit{IEEE TPAMI}, 14(2), 239--256.

\bibitem{rusu2009}
Rusu, R.~B., Blodow, N. \& Beetz, M. (2009).
Fast point feature histograms (FPFH) for 3D registration.
\textit{IEEE ICRA}, 3212--3217.

\bibitem{zhou2016}
Zhou, Q.-Y., Park, J. \& Koltun, V. (2016).
Fast global registration.
\textit{ECCV}, 766--782.

\bibitem{chen1992}
Chen, Y. \& Medioni, G. (1992).
Object modelling by registration of multiple range images.
\textit{Image and Vision Computing}, 10(3), 145--155.

\bibitem{open3d2018}
Zhou, Q.-Y., Park, J. \& Koltun, V. (2018).
Open3D: A modern library for 3D data processing.
\textit{arXiv:1801.09847}.
\end{thebibliography}

\vfill
\noindent\rule{\textwidth}{0.4pt}\\
\small\textit{Ipari Képfeldolgozás Lab.\ Gyak.\ | 2.~Mérföldkő | 2026.~április}

\end{document}
